\documentclass[12pt]{article}
\usepackage[utf8]{inputenc}
\usepackage[T1]{fontenc}
\usepackage[margin=1in]{geometry}
\usepackage{setspace}
\usepackage{xcolor}
\usepackage[colorlinks=true, urlcolor=blue!60!black, linkcolor=blue!60!black]{hyperref}
\setlength{\parindent}{2em}
\setlength{\parskip}{1.25em}
\renewcommand{\baselinestretch}{1.0}
\begin{document}
\begin{center}
{\LARGE \textit{Lochner v. New York}, 198 U.S. 45 (1905)\par}
\vspace{0.8em}
{\large Case Brief\par}
\end{center}
\vspace{1em}

\section{Memory Jogger}

In \href{https://www.courtlistener.com/opinion/96276/lochner-v-new-york/}{\textit{Lochner v. New York}}, the Supreme Court held that a New York law capping bakers' working hours at ten per day and sixty per week violated the ``liberty of contract'' protected by the Fourteenth Amendment's Due Process Clause because the statute lacked a direct and reasonable relation to public health or worker safety.

\section{Detailed Case Facts}

In 1895, and subsequently codified as Article~8, Chapter~415 of the Laws of 1897 (known as the New York Bakeshop Act), the New York State Legislature enacted comprehensive regulations governing commercial bakeries and confectionery establishments (\href{https://caselaw.findlaw.com/court/us-supreme-court/198/45.html}{198 U.S. at 46--47}). Sections~111 through~115 addressed sanitary conditions directly---requiring proper drainage, plumbing, ventilation, eight-foot ceilings, washable walls, separate washrooms and sleeping quarters, and periodic state inspection---while Section~110 imposed a strict maximum-hours cap: ``No employee shall be required or permitted to work in a biscuit, bread, or cake bakery or confectionery establishment more than sixty hours in any one week, or more than ten hours in any one day, unless for the purpose of making a shorter work day on the last day of the week'' (\href{https://caselaw.findlaw.com/court/us-supreme-court/198/45.html}{198 U.S. at 46 n.1}). Notably, Section~110 contained no exception for emergencies.

Joseph Lochner owned and operated a small retail bakery in Utica, New York (\href{https://civics.supremecourthistory.org/article/lochner-v-new-york/}{Supreme Court Historical Society}). Having previously been convicted of a first violation of the Bakeshop Act and fined \$25 (\href{https://www.oyez.org/cases/1900-1940/198us45}{Oyez}), Lochner was arrested again in April 1901 and indicted in Oneida County for ``wrongfully and unlawfully'' requiring and permitting an employee, Aman Schmitter, to work more than sixty hours in a single week (\href{https://caselaw.findlaw.com/court/us-supreme-court/198/45.html}{198 U.S. at 46--47}; \href{https://civics.supremecourthistory.org/article/lochner-v-new-york/}{Supreme Court Historical Society}). Lochner did not use physical compulsion; rather, as the state courts and the Supreme Court treated the record, the employee voluntarily agreed to work extra hours to earn additional wages (\href{https://caselaw.findlaw.com/court/us-supreme-court/198/45.html}{198 U.S. at 52--53}).

\section{Procedural History}

Lochner demurred to the indictment in the Oneida County Court, arguing that the facts stated did not constitute a crime because Section~110 violated the Fourteenth Amendment (\href{https://caselaw.findlaw.com/court/us-supreme-court/198/45.html}{198 U.S. at 47}). The trial court overruled the demurrer; when Lochner declined to plead further, the court entered a plea of not guilty, conducted a trial, convicted him of a misdemeanor (second offense), and sentenced him to pay a \$50 fine or stand committed to the Oneida County jail for up to fifty days (\href{https://caselaw.findlaw.com/court/us-supreme-court/198/45.html}{198 U.S. at 47}).

On appeal, the Appellate Division of the New York Supreme Court, Fourth Department, affirmed the conviction by a 3--2 vote, 73 App.~Div. 120 (1902), and the New York Court of Appeals affirmed by a closely divided 4--3 vote, 177 N.Y. 145 (1904), sustaining Section~110 as a valid health regulation under the state's police power (\href{https://caselaw.findlaw.com/court/us-supreme-court/198/45.html}{198 U.S. at 47, 58}). Lochner sought review in the U.S. Supreme Court by writ of error. The case was argued on February 23--24, 1905, and decided on April 17, 1905 (\href{https://www.oyez.org/cases/1900-1940/198us45}{Oyez}).

\section{Judicial Votes}

\textbf{Vote:} 5--4 reversing the New York Court of Appeals and invalidating Section~110 of the New York Labor Law (\href{https://www.oyez.org/cases/1900-1940/198us45}{Oyez}).

\begin{itemize}
    \item \textbf{Majority (5):} Justice Rufus W. Peckham delivered the opinion of the Court, joined by Chief Justice Melville W. Fuller and Justices David J. Brewer, Henry B. Brown, and Joseph McKenna.
    \item \textbf{Concurrences (0):} None.
    \item \textbf{Dissents (4):} Justice John Marshall Harlan filed a dissenting opinion, joined by Justices Edward D. White and William R. Day. Justice Oliver Wendell Holmes, Jr. filed a separate dissenting opinion.
\end{itemize}

\section{Holding}

\textbf{Judgment:} Reversed and remanded to the Oneida County Court for further proceedings not inconsistent with the Court's opinion (\href{https://caselaw.findlaw.com/court/us-supreme-court/198/45.html}{198 U.S. at 64--65}).

\textbf{Rule Announced:} The Due Process Clause of the Fourteenth Amendment protects the general ``liberty of contract''---including the reciprocal right of an employer to purchase labor and an employee to sell labor---as part of individual liberty (\href{https://caselaw.findlaw.com/court/us-supreme-court/198/45.html}{198 U.S. at 53, 56}). Although a state may curtail contractual freedom through a legitimate exercise of its police power (promoting public safety, health, morals, or general welfare), a regulation interfering with labor contracts between competent adults is unconstitutional unless it bears a ``direct relation, as a means to an end,'' to a legitimate police-power objective and rests on a fair, reasonable ground showing material danger to public or worker health (\href{https://caselaw.findlaw.com/court/us-supreme-court/198/45.html}{198 U.S. at 57--58, 61}).

\textbf{Holding vs.\ Dicta:} The narrow holding is that New York's sixty-hour weekly and ten-hour daily cap on bakery employees' working hours, unaccompanied by an emergency exception, is not a valid health regulation under the state police power and therefore violates the Fourteenth Amendment's Due Process Clause (\href{https://caselaw.findlaw.com/court/us-supreme-court/198/45.html}{198 U.S. at 64}). By contrast, Justice Peckham's observations endorsing the sanity and probable validity of the Bakeshop Act's separate sanitary provisions (Sections~111--115) (\href{https://caselaw.findlaw.com/court/us-supreme-court/198/45.html}{198 U.S. at 61--62}), as well as his cautionary hypotheticals warning that upholding Section~110 would empower legislatures to cap the hours of bank clerks, lawyers' clerks, doctors, and athletes (\href{https://caselaw.findlaw.com/court/us-supreme-court/198/45.html}{198 U.S. at 59--61}), are dicta illustrating the majority's slippery-slope rationale rather than binding rulings on unchallenged statutes.

\section{Analysis of Opinions}

\subsection{Majority Opinion (Justice Peckham)}

Justice Peckham began by characterizing Section~110 not as a rule against coerced labor, but as an absolute prohibition preventing willing employers and employees from contracting for overtime work (\href{https://caselaw.findlaw.com/court/us-supreme-court/198/45.html}{198 U.S. at 52--53}). Relying on \textit{Allgeyer v. Louisiana}, 165 U.S. 578 (1897), he declared that the general right to make contracts in relation to one's business---specifically including the right to buy and sell labor---is part of the ``liberty'' safeguarded by the Fourteenth Amendment (\href{https://caselaw.findlaw.com/court/us-supreme-court/198/45.html}{198 U.S. at 53}).

Peckham acknowledged that liberty of contract is not absolute and must yield to valid exercises of the state police power. He distinguished prior decisions upholding state regulations: \textit{Holden v. Hardy}, 169 U.S. 366 (1898), upheld an eight-hour day for underground miners and smelter workers because those occupations posed acute, peculiar health hazards and the Utah law included an emergency exception; \textit{Atkin v. Kansas}, 191 U.S. 207 (1903), involved a state's proprietary control over public works; and \textit{Jacobson v. Massachusetts}, 197 U.S. 11 (1905), upheld compulsory vaccination to protect the community against contagious disease (\href{https://caselaw.findlaw.com/court/us-supreme-court/198/45.html}{198 U.S. at 54--56}).

Applying active means-ends scrutiny, Peckham insisted that courts must independently verify whether a law is a ``fair, reasonable, and appropriate exercise of the police power'' or an ``unreasonable, unnecessary, and arbitrary interference'' with individual liberty; otherwise, the police power would become a ``mere pretext'' erasing Fourteenth Amendment limits (\href{https://caselaw.findlaw.com/court/us-supreme-court/198/45.html}{198 U.S. at 56}). First, as a ``labor law, pure and simple,'' Section~110 failed because bakers are competent adults ``in no sense wards of the state'' (\href{https://caselaw.findlaw.com/court/us-supreme-court/198/45.html}{198 U.S. at 57}). Second, as a public-health law, ``[c]lean and wholesome bread does not depend upon whether the baker works but ten hours per day or only sixty hours a week'' (\href{https://caselaw.findlaw.com/court/us-supreme-court/198/45.html}{198 U.S. at 57}). Third, as a worker-health measure, baking was not to ``common understanding'' an unusually unhealthy trade, and allowing legislatures to regulate hours in any job with minor health risks would place every occupation---from printers and carpenters to lawyers and bank clerks---``at the mercy of legislative majorities'' (\href{https://caselaw.findlaw.com/court/us-supreme-court/198/45.html}{198 U.S. at 59--60}). Pointing to the Act's separate sanitary sections (Sections~111--115) as proof that New York had already addressed bakery cleanliness directly, Peckham concluded that Section~110 was really passed ``from other motives''---namely, to regulate labor relations in a private business---and was therefore void (\href{https://caselaw.findlaw.com/court/us-supreme-court/198/45.html}{198 U.S. at 61--64}).

\subsection{Dissenting Opinion (Justice Harlan)}

Justice Harlan, joined by Justices White and Day, accepted the majority's doctrinal premise that the Fourteenth Amendment protects liberty of contract subject to reasonable police-power regulations (\href{https://caselaw.findlaw.com/court/us-supreme-court/198/45.html}{198 U.S. at 65--68}). He diverged sharply, however, on the standard of judicial review and the burden of proof. Invoking \textit{McCulloch v. Maryland}, \textit{Mugler v. Kansas}, and \textit{Jacobson v. Massachusetts}, Harlan argued that statutes enjoy a strong presumption of constitutionality and may be invalidated only if they have ``no real or substantial relation'' to public health, morals, or safety, or are ``beyond question, plainly and palpably in excess of legislative power'' (\href{https://caselaw.findlaw.com/court/us-supreme-court/198/45.html}{198 U.S. at 68}).

Marshaling medical treatises (including Professor Ludwig Hirt's \textit{Diseases of the Workers}), the Eighteenth Annual Report of the New York Bureau of Statistics of Labor, and comparative international labor data, Harlan documented that bakers suffered from inhaling flour dust, working long night shifts in overheated rooms, respiratory inflammation, and shortened life expectancy (\href{https://caselaw.findlaw.com/court/us-supreme-court/198/45.html}{198 U.S. at 70--72}). Because there was ample ``room for debate and for an honest difference of opinion'' over whether working more than sixty hours a week in a bakery endangered worker health, courts had a constitutional duty to defer to New York's legislative judgment (\href{https://caselaw.findlaw.com/court/us-supreme-court/198/45.html}{198 U.S. at 72--74}).

\subsection{Dissenting Opinion (Justice Holmes)}

Writing alone in one of the most famous dissents in American legal history, Justice Holmes attacked the majority's methodology at its root. He accused the Court of deciding the case ``upon an economic theory which a large part of the country does not entertain'' (\href{https://caselaw.findlaw.com/court/us-supreme-court/198/45.html}{198 U.S. at 75}). Reviewing accepted interferences with contractual freedom---Sunday laws, usury laws, lottery bans, compulsory schooling, vaccination mandates, antitrust restrictions, and stock-margin prohibitions---Holmes declared that ``[t]he Fourteenth Amendment does not enact Mr. Herbert Spencer's \textit{Social Statics}'' (\href{https://supreme.justia.com/cases/federal/us/198/45/}{198 U.S. at 75}).

Holmes articulated a sweeping philosophy of judicial restraint: a Constitution ``is not intended to embody a particular economic theory, whether of paternalism and the organic relation of the citizen to the state or of \textit{laissez faire}.'' Rather, ``[i]t is made for people of fundamentally differing views'' (\href{https://caselaw.findlaw.com/court/us-supreme-court/198/45.html}{198 U.S. at 75--76}). In Holmes's formulation, the word ``liberty'' in the Fourteenth Amendment is perverted when invoked to block ``the natural outcome of a dominant opinion, unless it can be said that a rational and fair man necessarily would admit that the statute proposed would infringe fundamental principles as they have been understood by the traditions of our people and our law'' (\href{https://caselaw.findlaw.com/court/us-supreme-court/198/45.html}{198 U.S. at 76}). Because a reasonable legislator could view Section~110 as either a health measure or a permissible first installment of general hours regulation, Holmes would have upheld the statute.

\section{Examples: Future Applications}

Applying \textit{Lochner}'s framework (assuming it remained controlling law during the \textit{Lochner} era) yields the following outcomes:

\begin{enumerate}
    \item \textbf{Same-Side (Invalidated): State Maximum-Hours Cap for Retail Store Clerks.} A state enacts a law prohibiting grocery and department stores from permitting adult clerks to work more than 48 hours per week, citing general fatigue and family well-being. Under \textit{Lochner}, this law would be struck down because retail clerking is an ordinary, non-hazardous occupation, adult clerks are \textit{sui juris}, and a general welfare or fatigue rationale lacks the direct, material health nexus required to override liberty of contract.

    \item \textbf{Same-Side (Invalidated): Mandatory Minimum Wage for Adult Restaurant Workers.} A state establishes a minimum hourly wage for adult servers and cooks in private restaurants to equalize bargaining power between employers and workers. Under \textit{Lochner}'s logic (later applied in \textit{Adkins v. Children's Hospital}, 261 U.S. 525 (1923), before being overruled in \textit{West Coast Hotel Co. v. Parrish}, 300 U.S. 379 (1937)), fixing wages in a private business is a ``pure'' labor regulation unrelated to public health or safety and therefore an unconstitutional interference with the right to buy and sell labor.

    \item \textbf{Opposite-Side (Upheld): Eight-Hour Daily Limit for Underground Chemical-Mine Workers with an Emergency Exception.} A state limits employees working in underground sulfur mines to eight hours per day, with an explicit exception for emergencies threatening life or property. Following \textit{Holden v. Hardy} as expressly distinguished by Justice Peckham in \textit{Lochner}, this law would be upheld because the severe, documented toxic hazards of underground mining establish a direct and reasonable connection between limiting hours and protecting worker health.

    \item \textbf{Opposite-Side (Upheld): Mandatory Ventilation, Sanitation, and Fire-Exit Standards for Commercial Bakeries.} A state requires all commercial bakeries to install mechanical exhaust fans, washable tile floors, and unobstructed fire exits, enforceable by factory inspectors. Even under \textit{Lochner}, Justice Peckham expressly indicated that direct sanitary and safety regulations of bakery premises (like Sections~111--115 of the Bakeshop Act) fall squarely within the state's legitimate police power because they directly protect public and worker health without dictating private wage-and-hour bargains.

    \item \textbf{Fence-Sitter (Close Call): Ten-Hour Daily Cap for Commercial Long-Haul Bus Drivers.} A state prohibits private intercity bus companies from allowing drivers to operate buses for more than ten hours per day. On one hand, a challenger under \textit{Lochner} would argue that driving is a common occupation performed by competent adults and that capping hours intrudes on labor contracts. On the other hand, unlike baking---where Justice Peckham saw no link between a baker's fatigue and the wholesomeness of bread---a fatigued bus driver poses an immediate, direct physical danger to passengers and third parties on public highways, giving the state a far stronger public-safety justification under the police power.
\end{enumerate}

\section{Critique}

\textit{Lochner v. New York} occupies a central place in the American constitutional ``anti-canon,'' yet legal scholars across ideological traditions disagree sharply over \textit{why} it was wrong---or whether parts of its reasoning deserve rehabilitation.

\textbf{Progressive and Doctrinal Critiques.} From the Progressive Era through the New Deal, critics echoed Justices Holmes and Harlan, condemning \textit{Lochner} as lawless judicial activism that substituted judges' \textit{laissez-faire} policy preferences for democratic legislation (\href{https://www.heritage.org/the-constitution/report/economic-liberty-and-the-constitution-introduction}{Heritage Foundation}; \href{https://supreme.justia.com/cases/federal/us/198/45/}{Justia}). In his influential 1987 essay, \href{https://chicagounbound.uchicago.edu/journal_articles/8453/}{\textit{Lochner's Legacy}}, Cass R. Sunstein refined this critique: the Court's core error was not merely ``activism,'' but treating common-law property and contract baselines as natural and neutral, while viewing statutory departures from those baselines as impermissible ``partisan'' redistribution (\href{https://en.wikipedia.org/wiki/Lochner_era}{Sunstein summary}). Constitutional historian Howard Gillman similarly argued that the \textit{Lochner}-era Court was enforcing a nineteenth-century anti-``class legislation'' tradition forbidding government favoritism toward particular economic groups---a framework that became unworkable in an industrialized economy marked by stark bargaining-power disparities (\href{https://en.wikipedia.org/wiki/Lochner_era}{Gillman summary}). The Supreme Court formally repudiated \textit{Lochner}'s economic substantive due process in \textit{Nebbia v. New York}, 291 U.S. 502 (1934), and \textit{West Coast Hotel Co. v. Parrish}, 300 U.S. 379 (1937), ushering in deferential rational-basis review for economic regulation (\href{https://civics.supremecourthistory.org/article/lochner-v-new-york/}{Supreme Court Historical Society}; \href{https://supreme.justia.com/cases/federal/us/198/45/}{Justia}).

\textbf{Originalist, Textualist, and Libertarian Perspectives.} Conservative originalists such as Robert Bork and Justice Antonin Scalia have also condemned \textit{Lochner}, but on textualist grounds: in their view, ``substantive due process'' is an oxymoron because the Due Process Clause guarantees fair \textit{procedures}, not unenumerated substantive rights, making \textit{Lochner} the doctrinal ancestor of later controversial substantive-due-process rulings like \textit{Roe v. Wade} (\href{https://chicagounbound.uchicago.edu/cgi/viewcontent.cgi?article=1022&context=uclrev_online}{Chicago Unbound}; \href{https://texaslawreview.org/in-defense-of-substantive-due-process/}{Texas Law Review}). By contrast, libertarian and revisionist scholars---most notably David E. Bernstein and Randy E. Barnett---challenge the conventional Progressive narrative. Bernstein argues that the \textit{Lochner} majority was not motivated by Social Darwinism or pro-corporate bias, but by a principled natural-rights and precedent-grounded commitment to limiting arbitrary police-power intrusions (noting that the Bakeshop Act itself benefited unionized large bakeries at the expense of immigrant-owned small shops) (\href{https://www.heritage.org/the-constitution/report/economic-liberty-and-the-constitution-introduction}{Heritage Foundation}; \href{https://www.researchgate.net/publication/228158799_Lochner_v_New_York_A_Centennial_Retrospective}{Bernstein}). Similarly, Barnett and Evan Bernick contend that while economic liberty may fit more naturally under the Privileges or Immunities Clause, the original spirit of due process properly bars arbitrary legislation resting on mere legislative will rather than constitutionally proper ends (\href{https://wmlawreview.org/sites/default/files/Barnett\%20\%26\%20Bernick.pdf}{Barnett \& Bernick}).

\textbf{Underlying Competing Values.} At bottom, \textit{Lochner} exposes an enduring constitutional tension between individual economic autonomy and democratic self-governance: how aggressively unelected federal judges should police the boundary between genuine public-health or safety measures and rent-seeking or paternalistic regulations, and whether the Constitution empowers courts to protect unenumerated economic liberties without usurping the legislature's authority to address industrial inequality.

\section{Key Quotations}

\begin{quote}
``Clean and wholesome bread does not depend upon whether the baker works but ten hours per day or only sixty hours a week. The limitation of the hours of labor does not come within the police power on that ground.''
\end{quote}
\noindent --- Justice Rufus W. Peckham, Opinion of the Court, \textit{Lochner v. New York}, 198 U.S. 45, 57 (1905) (\href{https://caselaw.findlaw.com/court/us-supreme-court/198/45.html}{FindLaw}; \href{https://www.courtlistener.com/opinion/96276/lochner-v-new-york/}{CourtListener}).

\begin{quote}
``The mere assertion that the subject relates, though but in a remote degree, to the public health, does not necessarily render the enactment valid. The act must have a more direct relation, as a means to an end, and the end itself must be appropriate and legitimate, before an act can be held to be valid which interferes with the general right of an individual to be free in his person and in his power to contract in relation to his own labor.''
\end{quote}
\noindent --- Justice Rufus W. Peckham, Opinion of the Court, \textit{Lochner v. New York}, 198 U.S. 45, 57--58 (1905) (\href{https://caselaw.findlaw.com/court/us-supreme-court/198/45.html}{FindLaw}).

\begin{quote}
``If there be doubt as to the validity of the statute, that doubt must therefore be resolved in favor of its validity, and the courts must keep their hands off, leaving the legislature to meet the responsibility for unwise legislation. If the end which the legislature seeks to accomplish be one to which its power extends, and if the means employed to that end, although not the wisest or best, are yet not plainly and palpably unauthorized by law, then the court cannot interfere.''
\end{quote}
\noindent --- Justice John Marshall Harlan, Dissenting, \textit{Lochner v. New York}, 198 U.S. 45, 68 (1905) (\href{https://caselaw.findlaw.com/court/us-supreme-court/198/45.html}{FindLaw}).

\begin{quote}
``The Fourteenth Amendment does not enact Mr. Herbert Spencer's \textit{Social Statics}. ... But a constitution is not intended to embody a particular economic theory, whether of paternalism and the organic relation of the citizen to the state or of \textit{laissez faire}. It is made for people of fundamentally differing views, and the accident of our finding certain opinions natural and familiar or novel and even shocking ought not to conclude our judgment upon the question whether statutes embodying them conflict with the Constitution of the United States.''
\end{quote}
\noindent --- Justice Oliver Wendell Holmes, Jr., Dissenting, \textit{Lochner v. New York}, 198 U.S. 45, 75--76 (1905) (\href{https://supreme.justia.com/cases/federal/us/198/45/}{Justia}; \href{https://caselaw.findlaw.com/court/us-supreme-court/198/45.html}{FindLaw}).

\begin{quote}
``General propositions do not decide concrete cases. The decision will depend on a judgment or intuition more subtle than any articulate major premise.''
\end{quote}
\noindent --- Justice Oliver Wendell Holmes, Jr., Dissenting, \textit{Lochner v. New York}, 198 U.S. 45, 76 (1905) (\href{https://caselaw.findlaw.com/court/us-supreme-court/198/45.html}{FindLaw}).

\end{document}
